\section{Related Work}

\subsection{Chatbots in Educational Settings}
Research on chatbots in educational contexts shows a shift from rule-based and highly structured conversational agents toward data-driven and increasingly knowledge-integrating systems \autocite{Waldner2022, Chen2025,Klesel2025}. Earlier work has focused primarily on rule-based or template-driven Pedagogical Conversational Agents (PCAs) used for structured tasks like quizzes \autocite{Waldner2022}. While offering high control, they lacked contextual adaptability and rarely promoted deeper learning \autocite{Waldner2022}. As data-driven models became widely available, the focus shifted toward tutor-like systems integrated into the learning process. Chatbots now adopt roles such as tutor \autocite{KhosrawiRad2024} or coach \autocite{Chen2022}, providing adaptive assistance and motivational interactions that significantly impact learning behavior \autocite{Schlimbach2023, KhosrawiRad2024}. To prevent cognitive offloading and deskilling, chatbots require a pedagogical redesign that actively supports cognitive processes rather than merely providing information \autocite{soc15010006}.

A key approach here is Socratic design, in which chatbots ask targeted questions to encourage reflection and independent thinking \autocite{socaraticChatbot}. As a counter-model to purely response-based systems, it reduces the risk of excessive reliance on AI-generated solutions \autocite{Ma2024}. Socratic agents are already used in commercial systems like Google's Gemini Tutor, DuoLingo's Max, and Khanmigo, with studies showing they enable personalized learning by scaffolding reasoning processes rather than delivering immediate answers \autocite{Shetye2024Khanmigo}. This transition from simple question-and-answer scenarios to interactive learning support places greater demands on system architectures, requiring context-sensitive retrieval and dynamic integration into dialogues.

These capabilities are increasingly implemented via Retrieval-Augmented Generation (RAG) \autocite{zhaoragSurvey}. RAG combines parametric language models with external knowledge to integrate domain-specific content \autocite{Lewis2020_RAG}. In education, course-specific chatbots use learning materials to generate contextualized responses, significantly increasing student interaction and improving perceived outcomes \autocite{Chen2025}. Thus, RAG is viewed as a foundational architecture for learning and support systems \autocite{Klesel2025,Chen2025}.

Recent research conceptualizes RAG beyond a single pipeline as a design space spanning retrieval, augmentation, and generation stages \autocite{zhaoragSurvey}. A distinction is made between naive, advanced, and modular systems. Naive RAG uses a linear pipeline to append retrieved documents to a prompt; its effectiveness heavily depends on retrieval quality. Advanced RAG improves this via hybrid retrieval, query rewriting, and reranking. Empirical evidence indicates retrieval quality often outweighs improvements from model scaling \autocite{liBPRAG}, making system-level design decisions—like chunking strategies—crucial for performance \autocite{zhaoragSurvey}.

Modular RAG dynamically orchestrates components for iterative retrieval and reasoning. This is particularly relevant for multi-hop reasoning, aggregating diverse sources through iterative queries \autocite{MultiHopRAG}. Further extensions like GraphRAG capture entity relationships as graphs to enable holistic reasoning, though introducing computational complexity \autocite{GraphRAG_Grundlage}.

\subsection{RAG System Evaluation in Educational Settings}

Evaluating RAG systems requires assessing retrieval quality, generation quality, and their interaction. Early evaluations relied on static QA benchmarks, which poorly reflect real-world deployments with continuous, domain-specific knowledge bases \autocite{Lewis2020_RAG}.

Recent research models RAG evaluation as a multi-stage problem. Retrieval quality focuses on metrics like Precision@k, Recall@k, and Mean Reciprocal Rank \autocite{manning2008ir, karpukhin2020dpr}. Generation quality addresses relevance and factual consistency. A central criterion here is \textit{faithfulness}, capturing whether outputs are genuinely backed by retrieved evidence \autocite{es2025ragas}. Explicit evaluation of factual consistency remains vital, as hallucinations still occur despite external references.

Furthermore, static, single-turn evaluations fail to capture multi-turn interaction complexities \autocite{Deriu2021}. Modern benchmarks like RGB test multi-document integration, noise robustness, and abstention behavior when evidence is lacking \autocite{chen2023benchmarkinglargelanguagemodels}, while CRUD-RAG evaluates components across various systemic tasks \autocite{CRUDRAG}.

To increase scalability, automatic evaluation frameworks using large language models (LLMs) as judges (e.g., RAGAS, ARES) estimate relevance and faithfulness without costly annotations \autocite{es2025ragas, saadfalcon2024ares}. However, due to prompt sensitivities and systematic biases, LLM-based evaluation requires caution as a standalone method. Hybrid approaches combining automated estimation with expert human judgment are highly recommended to balance scalability with nuanced qualitative insights \autocite{bdcc9120320, 202504.0369}.

In broader educational contexts, evaluation merges technical correctness with pedagogical quality and user-centered metrics like perceived usefulness \autocite{Neumann2024MoodleBot}. Consequently, applied evaluations of RAG-based systems synthesize technical benchmarks (faithfulness, relevance, latency) with user-facing dimensions \autocite{khan2025ragbenchmark}.

Overall, the literature underscores that multidimensional evaluation is crucial. However, studies still heavily rely on standardized rather than realistic, application-oriented datasets. Ultimately, cross-architectural evaluations that systematically balance automation, human expertise, and operational costs (such as latency) are scarce. Addressing these gaps directly motivates the framework proposed in this work.